\subsubsection{Penilaian Ancaman dengan DREAD}

Ancaman yang telah diklasifikasikan pada Tabel~\ref{tab:klasifikasi-ancaman} selanjutnya dinilai tingkat risikonya menggunakan metode DREAD (\textit{Damage Potential}, \textit{Reproducibility}, \textit{Exploitability}, \textit{Affected Users}, \textit{Discoverability}). Setiap kategori diberi bobot Tinggi~=~3, Sedang~=~2, atau Rendah~=~1 sesuai Tabel~\ref{tab:dread-rating-reference}.

Berikut adalah tabel referensi penilaian yang akan digunakan untuk penilaian para ancaman ke kategori DREAD.

\begin{xltabular}{\textwidth}{|>{\centering\arraybackslash}p{0.6cm}|>{\raggedright\arraybackslash}p{2.2cm}|>{\raggedright\arraybackslash}X|>{\raggedright\arraybackslash}X|>{\raggedright\arraybackslash}X|}
    \caption{Tabel Referensi Penilaian DREAD}
    \label{tab:dread-rating-reference} \\
    \hline
    & \textbf{Rating} &
    \textbf{Tinggi (3)} &
    \textbf{Sedang (2)} &
    \textbf{Rendah (1)} \\
    \hline
    \endfirsthead

    \multicolumn{5}{c}%
    {} \\
    \hline
    & \textbf{Rating} &
    \textbf{Tinggi (3)} &
    \textbf{Sedang (2)} &
    \textbf{Rendah (1)} \\
    \hline
    \endhead

    \hline
    \multicolumn{5}{r}{} \\
    \endfoot

    \hline
    \endlastfoot

    D &
    Damage Potential &
    Penyerang dapat menumbangkan sistem keamanan; memperoleh otorisasi kepercayaan penuh; berjalan sebagai administrator; mengunggah konten. &
    Kebocoran informasi yang sensitif. &
    Kebocoran informasi yang tidak signifikan (trivial). \\
    \hline

    R &
    Reproduci\-bility &
    Serangan dapat direproduksi setiap saat dan tidak memerlukan jendela waktu (timing window) tertentu. &
    Serangan dapat direproduksi, tetapi hanya dengan jendela waktu dan kondisi race tertentu. &
    Serangan sangat sulit direproduksi, bahkan dengan pengetahuan atas celah keamanannya. \\
    \hline

    E &
    Exploitability &
    Seorang programmer pemula dapat melakukan serangan ini dalam waktu singkat. &
    Seorang programmer terampil dapat melakukan serangan ini, kemudian mengulangi langkah-langkahnya. &
    Serangan memerlukan orang yang sangat terampil dan pengetahuan mendalam setiap kali akan dieksploitasi. \\
    \hline

    A &
    Affected Users &
    Seluruh pengguna, konfigurasi default, pelanggan utama (key customers). &
    Sebagian pengguna, konfigurasi non-default. &
    Persentase pengguna yang sangat kecil, fitur yang jarang digunakan; memengaruhi pengguna anonim. \\
    \hline

    D &
    Discover - ability &
    Informasi mengenai serangan telah dipublikasikan. Kerentanan ditemukan pada fitur yang paling umum digunakan dan sangat mudah terlihat. &
    Kerentanan berada pada bagian produk yang jarang digunakan, dan hanya sedikit pengguna yang mungkin menemukannya. Diperlukan pemikiran lebih untuk melihat potensi penyalahgunaannya. &
    Bug bersifat tersembunyi (obscure), dan pengguna kemungkinan besar tidak akan menyadari potensi kerusakannya. \\
    \hline

\end{xltabular}

% Sumber: Howard, M.; LeBlanc, D. Improving Web Application Security: Threats and Countermeasures; Microsoft Press, 2003.

Beberapa ancaman melibatkan lebih dari satu level kepercayaan yang berpotensi menghasilkan nilai \textit{Damage Potential} dan/atau \textit{Affected Users} yang berbeda -- misalnya ancaman \textit{spoofing} terhadap kredensial \textit{sign-in} akan menimbulkan kerusakan yang lebih besar apabila menimpa Ministry Administrator (LK6) dibandingkan Pasien (LK2). Pada kondisi tersebut, ancaman dinilai secara terpisah untuk setiap level kepercayaan yang menghasilkan nilai DREAD berbeda. Sebaliknya, apabila nilai DREAD tidak berbeda antar level kepercayaan, penilaian digabungkan dalam satu baris.

Nilai total (\textbf{JML}) merupakan penjumlahan kelima kategori dengan rentang 5--15, yang kemudian dipetakan menjadi peringkat risiko sesuai Tabel~\ref{tab:peringkat-risiko}.

\begin{table}[h]
\centering
\caption{Peringkat Risiko DREAD}
\label{tab:peringkat-risiko}
\begin{tabular}{|c|c|}
\hline
\textbf{Nilai} & \textbf{Peringkat Risiko} \\
\hline
5 -- 7 & Rendah \\
\hline
8 -- 11 & Sedang \\
\hline
12 -- 15 & Tinggi \\
\hline
\end{tabular}
\end{table}

Hasil penilaian ancaman selengkapnya disajikan pada Tabel~\ref{tab:penilaian-ancaman}. Kolom \textbf{D} pertama merepresentasikan \textit{Damage Potential} dan kolom \textbf{D} kedua merepresentasikan \textit{Discoverability}, mengikuti urutan akronim DREAD.

\begin{xltabular}{\textwidth}{|>{\centering\arraybackslash}p{1cm}|>{\raggedright\arraybackslash}p{3.1cm}|>{\centering\arraybackslash}p{0.5cm}|>{\centering\arraybackslash}p{0.5cm}|>{\centering\arraybackslash}p{0.5cm}|>{\centering\arraybackslash}p{0.5cm}|>{\centering\arraybackslash}p{0.5cm}|>{\centering\arraybackslash}p{0.9cm}|>{\centering\arraybackslash}p{1.6cm}|}
    \caption{Penilaian Ancaman Sistem DecMed} \label{tab:penilaian-ancaman} \\
    
    \hline
    \textbf{ID} & \textbf{Level Kepercayaan} & \textbf{D} & \textbf{R} & \textbf{E} & \textbf{A} & \textbf{D} & \textbf{JML} & \textbf{Risiko} \\
    \hline
    \endfirsthead
 
    \hline
    \textbf{ID} & \textbf{Level Kepercayaan} & \textbf{D} & \textbf{R} & \textbf{E} & \textbf{A} & \textbf{D} & \textbf{JML} & \textbf{Risiko} \\
    \hline
    \endhead
 
    \hline
    \endfoot
 
    \hline
    \endlastfoot
 
    T1 & LK2 (Pasien) & 1 & 3 & 3 & 1 & 2 & 10 & Sedang \\ \hline
    T1 & LK3 (Adm. Personnel) & 1 & 3 & 3 & 2 & 2 & 11 & Sedang \\ \hline
    T1 & LK4, LK5 (Medis, Hosp. Admin) & 2 & 3 & 3 & 2 & 2 & 12 & Tinggi \\ \hline
    T1 & LK6 (Ministry Admin) & 3 & 3 & 3 & 3 & 2 & 14 & Tinggi \\ \hline
 
    T2 & LK2 (Pasien) & 1 & 2 & 2 & 1 & 2 & 8 & Sedang \\ \hline
    T2 & LK3 (Adm. Personnel) & 1 & 2 & 2 & 2 & 2 & 9 & Sedang \\ \hline
    T2 & LK4, LK5 (Medis, Hosp. Admin) & 2 & 2 & 2 & 2 & 2 & 10 & Sedang \\ \hline
    T2 & LK6 (Ministry Admin) & 3 & 2 & 2 & 3 & 2 & 12 & Tinggi \\ \hline
 
    T3 & LK1 & 1 & 3 & 2 & 1 & 2 & 9 & Sedang \\ \hline
 
    T4 & LK2, LK4 & 2 & 2 & 1 & 1 & 1 & 7 & Rendah \\ \hline
 
    T5 & LK2, LK4 & 1 & 2 & 1 & 1 & 1 & 6 & Rendah \\ \hline
 
    T6 & LK2 & 1 & 2 & 1 & 1 & 1 & 6 & Rendah \\ \hline
 
    T7 & LK3 (Adm. Personnel) & 1 & 2 & 1 & 1 & 1 & 6 & Rendah \\ \hline
    T7 & LK4 (Medical Personnel) & 2 & 2 & 1 & 2 & 1 & 8 & Sedang \\ \hline
 
    T8 & LK5 & 2 & 2 & 1 & 2 & 1 & 8 & Sedang \\ \hline
 
    T9 & LK6 & 2 & 1 & 1 & 1 & 1 & 6 & Rendah \\ \hline
 
    T10 & LK2 (Pasien) & 2 & 1 & 1 & 1 & 1 & 6 & Rendah \\ \hline
    T10 & LK3, LK4 (Personnel) & 3 & 1 & 1 & 2 & 1 & 8 & Sedang \\ \hline
 
    T11 & LK2 (Pasien) & 1 & 2 & 1 & 1 & 1 & 6 & Rendah \\ \hline
    T11 & LK3, LK4 (Personnel) & 2 & 2 & 1 & 2 & 1 & 8 & Sedang \\ \hline
 
    T12 & LK1 & 2 & 3 & 2 & 2 & 2 & 11 & Sedang \\ \hline
 
    T13 & LK3 (Adm. Personnel) & 2 & 2 & 1 & 1 & 1 & 7 & Rendah \\ \hline
    T13 & LK4 (Medical Personnel) & 3 & 2 & 1 & 2 & 1 & 9 & Sedang \\ \hline
 
    T14 & LK7 & 2 & 1 & 1 & 1 & 1 & 6 & Rendah \\ \hline
 
    T15 & LK7, LK8 & 3 & 1 & 1 & 2 & 1 & 8 & Sedang \\ \hline
 
    T16 & LK7, LK8 & 2 & 1 & 1 & 2 & 1 & 7 & Rendah \\ \hline
 
    T17 & LK7, LK8 & 1 & 1 & 1 & 1 & 1 & 5 & Rendah \\ \hline
 
    T18 & LK1--LK4 & 2 & 2 & 2 & 2 & 2 & 10 & Sedang \\ \hline
 
    T19 & Kunci milik Pasien (LK2) & 2 & 1 & 1 & 1 & 1 & 6 & Rendah \\ \hline
    T19 & Kunci milik Tenaga Fasyankes (LK3, LK4) & 2 & 1 & 1 & 2 & 1 & 7 & Rendah \\ \hline
    T19 & Kunci milik Hospital Admin (LK5) & 3 & 1 & 1 & 2 & 1 & 8 & Sedang \\ \hline
 
    T20 & Kunci milik Pasien (LK2) & 1 & 1 & 1 & 1 & 1 & 5 & Rendah \\ \hline
    T20 & Kunci milik Tenaga Fasyankes (LK3, LK4) & 2 & 1 & 1 & 2 & 1 & 7 & Rendah \\ \hline
    T20 & Kunci milik Hospital Admin (LK5) & 3 & 1 & 1 & 2 & 1 & 8 & Sedang \\ \hline
 
    T21 & LK7 & 3 & 1 & 1 & 2 & 1 & 8 & Sedang \\ \hline
 
    T22 & LK7 & 2 & 1 & 1 & 2 & 1 & 7 & Rendah \\ \hline
 
    T23 & LK1 & 2 & 2 & 2 & 2 & 2 & 10 & Sedang \\ \hline
 
    T24 & LK1, publik & 2 & 2 & 1 & 2 & 2 & 9 & Sedang \\ \hline
 
    T25 & publik & 2 & 2 & 2 & 2 & 2 & 10 & Sedang \\ \hline
 
    T26 & LK9, publik & 1 & 3 & 2 & 2 & 3 & 11 & Sedang \\ \hline
 
    T27 & LK2/LK4 $\rightarrow$ LK6/LK7 & 3 & 1 & 1 & 3 & 1 & 9 & Sedang \\ \hline
 
    T28 & LK3 $\rightarrow$ LK4 & 3 & 2 & 1 & 2 & 1 & 9 & Sedang \\ \hline
 
\end{xltabular}
% 